\appendix

\chapter{Mã nguồn thực nghiệm}
\label{app:code}

Toàn bộ phần thực nghiệm của Chương~\ref{chap:thucnghiem} được cài đặt bằng
PyTorch thuần, không dùng thư viện PINNs đóng gói sẵn, và được công bố kèm báo
cáo trong thư mục \texttt{code/}. Phụ lục này mô tả cấu trúc và các quy ước cần
biết để tái lập kết quả.

%==============================================================================
\section{Cấu trúc và cách chạy}
%==============================================================================

\begin{table}[H]
\centering
\caption{Tổ chức thư mục \texttt{code/}.}
\label{tab:code-struct}
\small
\begin{tabular}{@{}l p{9.2cm}@{}}
\toprule
\textbf{Tệp} & \textbf{Nội dung} \\
\midrule
\texttt{pinns/core.py} & Thành phần dùng chung: lớp \texttt{FNN}, toán tử đạo
  hàm lồng nhau \texttt{grad}, chỉ tiêu sai số, chỉ số mất cân bằng
  $\rho_i$, và hai giai đoạn huấn luyện Adam~$\to$~L-BFGS. \\
\texttt{pinns/exp1\_relu.py} & Mục~\ref{sec:tn1} \\
\texttt{pinns/exp2\_lambda.py} & Mục~\ref{sec:tn2} \\
\texttt{pinns/exp3\_spectral.py} & Mục~\ref{sec:tn3} \\
\texttt{pinns/exp4\_heat.py} & Mục~\ref{sec:tn4} \\
\texttt{pinns/exp5\_burgers.py} & Mục~\ref{sec:tn5}, kèm bộ giải tham chiếu sai
  phân hữu hạn \\
\texttt{pinns/exp6\_stability.py} & Mục~\ref{sec:tn6} \\
\texttt{pinns/exp7\_nu\_sweep.py} & Mục~\ref{sec:tn7}, quét $\nu$ trên ba hạt
  giống \\
\texttt{pinns/exp8\_seeds.py} & Mục~\ref{sec:tn8}, lặp lại TN2, TN3, TN4, TN5
  trên năm hạt giống; chạy tiếp được sau khi bị ngắt, vì ghi kết quả sau mỗi hạt
  giống \\
\texttt{pinns/exp9\_chiphi.py} & Mục~\ref{sec:tn9}, đo chi phí một bước, tổng
  thời gian huấn luyện, và thời gian bộ giải sai phân hữu hạn ở cùng mức sai số \\
\texttt{pinns/exp10\_burgers\_manh.py} & Mục~\ref{sec:tn10}, PINN Burgers với
  mạng $8\times20$, $N_r = 10\,000$, L-BFGS nhiều vòng và RAR, trên năm hạt
  giống; nhận số hạt giống làm tham số để chạy song song, và \texttt{tong\_ket}
  để gộp kết quả và đo chi phí sai phân ở cùng mức sai số \\
\texttt{pinns/exp10b\_rar\_tachbien.py} & Mục~\ref{sec:tn10}, thí nghiệm tách
  biến: rẽ ba nhánh từ cùng trạng thái sau đợt L-BFGS thứ nhất --- RAR, thêm
  điểm ngẫu nhiên cùng số lượng, không thêm điểm \\
\texttt{pinns/exp10c\_rad.py} & Mục~\ref{sec:tn10}, biến thể lấy mẫu lại toàn
  bộ (RAD) với $c = 1$ và $c = 0$, rẽ từ cùng trạng thái như
  \texttt{exp10b} \\
\texttt{pinns/animate.py} & Hoạt hình quá trình huấn luyện (TN1, TN3, TN10) ra
  GIF và PDF nhiều trang cho slide, cùng Hình~\ref{fig:tn10} \\
\texttt{pinns/run\_all.py} & Chạy lần lượt cả mười thí nghiệm \\
\texttt{results/*.json} & Số liệu thô của từng thí nghiệm \\
\bottomrule
\end{tabular}
\end{table}

Mỗi thí nghiệm chạy độc lập bằng một lệnh dạng
\texttt{python -m pinns.<tệp>}, còn \texttt{python -m pinns.run\_all}
chạy toàn bộ. Kết quả số được ghi ra \texttt{results/} dưới dạng JSON để có thể
đối chiếu lại với các bảng trong báo cáo.

%==============================================================================
\section{Các quy ước bám theo phần lý thuyết}
%==============================================================================

Mọi lựa chọn thiết kế trong \texttt{core.py} đều truy được về một kết quả ở
Chương~\ref{chap:fnn} hoặc Chương~\ref{chap:autodiff}, đúng theo tinh thần đã nêu
ở Bảng~\ref{tab:tong-ket}: hàm kích hoạt $\tanh$ (điều kiện
$\sigma''\not\equiv 0$ của Hệ quả~\ref{hq:relu-hessian}); khởi tạo Xavier với
phương sai \eqref{eq:xavier}; độ chệch bằng $\bm{0}$ (giả thiết (A2) của
Mục~\ref{subsec:phuong-sai}); lớp ra tuyến tính \eqref{eq:fnn-output}; số học
\texttt{float64}; và huấn luyện hai giai đoạn theo
Thuật toán~\ref{alg:two-stage} với tìm kiếm đường Wolfe mạnh, vốn là điều kiện để
$\bm{H}_k$ giữ tính xác định dương (Bổ đề~\ref{bd:wolfe-curvature}).

\begin{nhanxet}[Quy ước lưới điểm phối trí]
\label{nx:quy-uoc-luoi}
Lưới điểm phối trí một chiều bao gồm \emph{cả hai điểm biên}. Chi tiết này không
tuỳ tiện và có thể kiểm chứng: với $n = 256$ trên $(0,1)$ và
$f = 4\pi^2\sin(2\pi x)$, trung bình của $f^2$ trên lưới ấy bằng đúng
$776{,}229$ --- trùng với giá trị $J_r$ mà Mục~\ref{sec:tn1} báo cáo cho mạng
ReLU. Nếu dùng $256$ điểm \emph{trong} thay vì kể cả biên, giá trị tương ứng là
$782{,}317$. Ta ghi rõ quy ước này vì nó là loại chi tiết quyết định việc một
bảng số có tái lập được hay không.
\end{nhanxet}

%==============================================================================
\section{Hai đoạn mã trung tâm}
%==============================================================================

Toán tử đạo hàm lồng nhau ở Đoạn mã~\ref{lst:grad} là hiện thực trực tiếp của hai
hệ thức truy hồi \eqref{eq:G-truy-hoi} và \eqref{eq:S-truy-hoi}; cờ
\texttt{create\_graph=True} là điều kiện để lấy tiếp gradient theo $\btheta$
(Mục~\ref{subsec:bac-cao-ad}). Lớp mạng ở Đoạn mã~\ref{lst:fnn} hiện thực
Định nghĩa~\ref{def:fnn} cùng ba quy ước vừa nêu.

Riêng với Mục~\ref{sec:tn6}, nghiệm thử ràng buộc cứng của
Bổ đề~\ref{bd:rang-buoc-cung} được dựng bằng đúng một dòng:

\begin{lstlisting}[caption={Nghiệm thử thoả điều kiện biên theo cấu trúc.},label={lst:hard}]
# u(0) = u(1) = 0 voi MOI theta, theo Bo de rang buoc cung
model = lambda x: x * (1 - x) * net(x)
\end{lstlisting}

%==============================================================================
\section{Giới hạn về tính tái lập}
%==============================================================================

Mọi thí nghiệm dùng số học \texttt{float64}; các mục không quét hạt giống dùng
\texttt{seed = 0}, còn Mục~\ref{sec:tn8} quét $\texttt{seed} \in \{0,\dots,4\}$
và Mục~\ref{sec:tn7} quét ba hạt giống. Mục~\ref{sec:tn8} lặp lại TN2, TN3, TN4
và TN5 trên $\texttt{seed} \in \{0,\dots,4\}$, tức mọi thí nghiệm của chương có
phụ thuộc vào khởi tạo hoặc lấy mẫu.

Cố định hạt giống \emph{vẫn chưa đủ}. Toàn bộ số liệu kèm theo được sinh với
\begin{center}
\texttt{OMP\_NUM\_THREADS=1 python -m pinns.<tệp>}
\end{center}
và biến môi trường ấy là bắt buộc chứ không phải tuỳ chọn: số luồng của thư viện
đại số tuyến tính thay đổi thứ tự cộng dồn trong phép nhân ma trận, và trên bài
toán Burgers điều đó đủ để đổi sai số cuối gần hai lần
(Mục~\ref{sec:hanche}, hạn chế~(7)).

Cần đọc cả hai điều trên kèm hạn chế thứ nhất ở Mục~\ref{sec:hanche}: tái lập
được không đồng nghĩa với có ý nghĩa thống kê, và $n = 5$ chỉ đủ để nêu trung vị
cùng khoảng chứ chưa đủ cho một khoảng tin cậy.
